\documentclass[10pt,a4paper]{amsart}
\usepackage{amsmath,amssymb,amsthm,graphicx}
% For commutative diagrams we use Xy-pic
\usepackage[all]{xy}

% We use 2cell for 2-commutative diagrams.
\xyoption{2cell}
% Theorem environments.
%
\theoremstyle{plain}
\newtheorem{theorem}[subsection]{Theorem}
\newtheorem{proposition}[subsection]{Proposition}
\newtheorem{lemma}[subsection]{Lemma}

\theoremstyle{definition}
\newtheorem{definition}[subsection]{Definition}
\newtheorem{example}[subsection]{Example}
\newtheorem{exercise}[subsection]{Exercise}
\newtheorem{situation}[subsection]{Situation}

\theoremstyle{remark}
\newtheorem{remark}[subsection]{Remark}
\newtheorem{remarks}[subsection]{Remarks}

\numberwithin{equation}{subsection}

% Macros
%
\def\lim{\mathop{\mathrm{lim}}\nolimits}
\def\colim{\mathop{\mathrm{colim}}\nolimits}
\def\Spec{\mathop{\mathrm{Spec}}}
\def\Hom{\mathop{\mathrm{Hom}}\nolimits}
\def\Ext{\mathop{\mathrm{Ext}}\nolimits}
\def\SheafHom{\mathop{\mathcal{H}\!\mathit{om}}\nolimits}
\def\SheafExt{\mathop{\mathcal{E}\!\mathit{xt}}\nolimits}
\def\Sch{\mathit{Sch}}
\def\Mor{\mathop{\mathrm{Mor}}\nolimits}
\def\Ob{\mathop{\mathrm{Ob}}\nolimits}
\def\Sh{\mathop{\mathit{Sh}}\nolimits}
\def\NL{\mathop{N\!L}\nolimits}
\def\CH{\mathop{\mathrm{CH}}\nolimits}
\def\proetale{{pro\text{-}\acute{e}tale}}
\def\etale{{\acute{e}tale}}
\def\QCoh{\mathit{QCoh}}
\def\Ker{\mathop{\mathrm{Ker}}}
\def\Im{\mathop{\mathrm{Im}}}
\def\Coker{\mathop{\mathrm{Coker}}}
\def\Coim{\mathop{\mathrm{Coim}}}

% Boxtimes
%
\DeclareMathSymbol{\boxtimes}{\mathbin}{AMSa}{"02}

%
% Macros for moduli stacks/spaces
%
\def\QCohstack{\mathcal{QC}\!\mathit{oh}}
\def\Cohstack{\mathcal{C}\!\mathit{oh}}
\def\Spacesstack{\mathcal{S}\!\mathit{paces}}
\def\Quotfunctor{\mathrm{Quot}}
\def\Hilbfunctor{\mathrm{Hilb}}
\def\Curvesstack{\mathcal{C}\!\mathit{urves}}
\def\Polarizedstack{\mathcal{P}\!\mathit{olarized}}
\def\Complexesstack{\mathcal{C}\!\mathit{omplexes}}
% \Pic is the operator that assigns to X its picard group, usage \Pic(X)
% \Picardstack_{X/B} denotes the Picard stack of X over B
% \Picardfunctor_{X/B} denotes the Picard functor of X over B
\def\Pic{\mathop{\mathrm{Pic}}\nolimits}
\def\Picardstack{\mathcal{P}\!\mathit{ic}}
\def\Picardfunctor{\mathrm{Pic}}
\def\Deformationcategory{\mathcal{D}\!\mathit{ef}}
\usepackage{cleveref}
\crefname{originalreference}{original reference}{original references}
\Crefname{originalreference}{Original reference}{Original references}
\makeatletter
\newlabel{lemma-directed-colimit-cohomology}{{03Q6}{source}{}{sourcepointer-1}{}}
\newlabel{lemma-directed-colimit-cohomology@cref}{{[originalreference][0][]03Q6}{[1][0][]source}{}{}{}}
\makeatother
\crefname{originalreference}{original reference}{original references}
\Crefname{originalreference}{Original reference}{Original references}
\makeatletter
\newlabel{lemma-gabber-for-absolutely-algebraically-closed}{{09ZC}{source}{}{sourcepointer-2}{}}
\newlabel{lemma-gabber-for-absolutely-algebraically-closed@cref}{{[originalreference][0][]09ZC}{[1][0][]source}{}{}{}}
\makeatother
\crefname{originallemma}{original lemma}{original lemmas}
\Crefname{originallemma}{Original lemma}{Original lemmas}
\makeatletter
\newlabel{limits-lemma-integral-limit-finite-and-finite-presentation}{{09YZ}{source}{}{sourcepointer-3}{}}
\newlabel{limits-lemma-integral-limit-finite-and-finite-presentation@cref}{{[originallemma][0][]09YZ}{[1][0][]source}{}{}{}}
\makeatother
\title{Affine schemes admit an integral cover killing all higher cohomology of finite constant coefficients, and any class is killed by a finite finitely presented cover}
\author{The Stacks Project Authors}
\date{stacks-project:etale-cohomology.tex:local-snapshot-20261011; GNU FDL 1.2 or later}
\begin{document}
\maketitle
\noindent\textit{Editorial note.} All delivered mathematical text is the author's own, copied from the pinned Stacks Project chapter file etale-cohomology.tex; no statement or proof step was rewritten, reordered or removed.\par\medskip
\noindent\textit{Editorial note.} The counted proof is the author's complete argument: the proof is written out in full: the construction of the integral cover through the integral closure of a polynomial ring modulo J, the appeal to the vanishing statement over absolutely algebraically closed rings, the presentation of the cover as a filtered limit of finite finitely presented morphisms, and the descent of the vanishing to a finite level.\par\medskip
\noindent\textit{Editorial note.} The label lemma-integral-cover-trivial-cohomology names the counted statement, and the proof environment immediately following it in the source is delivered verbatim as its proof.\par\medskip
\noindent\textit{Editorial note.} The counted statement is carried by the author's own lemma environment in the chapter source.\par\medskip
\noindent\textit{Editorial note.} The statement and proof are copied from the pinned chapter file; the three cited inputs are kept as external pointers with their canonical Stacks tags.\par\medskip
\noindent\textit{Editorial note.} The preamble is the chapter's own native preamble excerpt already used for this source: the Xy-pic setup (native preamble lines 14-18) and the native theorem declarations with the mathematical macros (native preamble lines 154-217).\par\medskip
\noindent\textit{Allowed standard prerequisites (editorial).}\begin{itemize}
\item Absolute integral closure and the vanishing statement over it. Same chapter, the lemma on cohomology of constant coefficients over schemes with absolutely algebraically closed rings. Supplies part (1) once the integral closure is set up.
\item Presentation of an integral morphism as a limit of finite finitely presented ones. Limits, the lemma presenting an integral morphism as a filtered limit of finite morphisms of finite presentation. Supplies X' = lim X'\_i in part (2).
\item Cohomology of directed limits. Same chapter, the lemma computing the cohomology of a directed limit of schemes as the colimit of the cohomologies. Supplies that a class vanishing on the limit already vanishes at some stage.
\end{itemize}\medskip
\ifdefined\hypertarget\hypertarget{sourcepointer-1}{}\fi
\noindent\textit{Source reference.} Original reference 03Q6: sources/stacks-etale-cohomology/pinned-original-source.tex, source lines 7464--7464. Prerequisite pointer invoked at proof line(s) 14495; it names a standard fact of the same chapter that the counted argument uses, and every step of the counted proof built on it is retained verbatim in the delivered body.\par\medskip
\ifdefined\hypertarget\hypertarget{sourcepointer-2}{}\fi
\noindent\textit{Source reference.} Original reference 09ZC: sources/stacks-etale-cohomology/pinned-original-source.tex, source lines 14445--14445. Prerequisite pointer invoked at proof line(s) 14485; it names a standard fact of the same chapter that the counted argument uses, and every step of the counted proof built on it is retained verbatim in the delivered body.\par\medskip
\ifdefined\hypertarget\hypertarget{sourcepointer-3}{}\fi
\noindent\textit{Source reference.} Original lemma 09YZ: sources/stacks-chow/tags.txt, source lines 12582--12582. Prerequisite pointer invoked at proof line(s) 14494; it names a standard fact of a sibling chapter that the counted argument uses, and every step of the counted proof built on it is retained verbatim in the delivered body.\par\medskip
\begin{lemma}
\label{lemma-integral-cover-trivial-cohomology}
Let $X$ be an affine scheme.
\begin{enumerate}
\item There exists an integral surjective morphism $X' \to X$ such that for
every closed subscheme $Z' \subset X'$, every finite abelian group $M$, and
every $q \geq 1$ we have $H^q_\etale(Z', \underline{M}) = 0$.
\item For any closed subscheme $Z \subset X$, finite abelian group $M$,
$q \geq 1$, and $\xi \in H^q_\etale(Z, \underline{M})$ there exists a
finite surjective morphism $X' \to X$ of finite presentation such that
$\xi$ pulls back to zero in $H^q_\etale(X' \times_X Z, \underline{M})$.
\end{enumerate}
\end{lemma}

\begin{proof}
Write $X = \Spec(A)$. Write $A = \mathbf{Z}[x_i]/J$ for some ideal $J$.
Let $R$ be the integral closure of $\mathbf{Z}[x_i]$ in an algebraic
closure of the fraction field of $\mathbf{Z}[x_i]$. Let
$A' = R/JR$ and set $X' = \Spec(A')$. This gives an example as in (1) by
Lemma \ref{lemma-gabber-for-absolutely-algebraically-closed}.

\medskip\noindent
Proof of (2). Let $X' \to X$ be the integral surjective morphism we found
above. Certainly, $\xi$ maps to zero in
$H^q_\etale(X' \times_X Z, \underline{M})$. We may write $X'$ as a
limit $X' = \lim X'_i$ of schemes finite and of finite presentation
over $X$; this is easy to do in our current affine case, but it
is a special case of the more general Limits, Lemma
\ref{limits-lemma-integral-limit-finite-and-finite-presentation}.
By Lemma \ref{lemma-directed-colimit-cohomology}
we see that $\xi$ maps to zero in $H^q_\etale(X'_i \times_X Z, \underline{M})$
for some $i$ large enough.
\end{proof}

\end{document}
